\documentclass[11pt]{article}
\usepackage[margin=1in]{geometry}
\usepackage{amsmath,amssymb,booktabs,url,xcolor}
\usepackage[hidelinks]{hyperref}
\usepackage{graphicx}
\emergencystretch=1.5em

\title{Staged Upgrades of Two Learned Populations: Recovering Part of a Learned Fleet's Benefit When the Signals It Relies On Are Upgraded\\[4pt]
\large DRAFT (05/10/2026)}
\date{}

\begin{document}
\maketitle

\begin{abstract}
Learned traffic signals and learned connected-vehicle fleets are deployed and upgraded on different schedules. In SUMO, a learned
eco-driving fleet that saves fuel under the old signal controller loses that benefit when the signal is upgraded to a learned controller.
The obvious repairs fail: withdrawing the timing returns the fleet to its base driver; forecasting the learned signal's timing costs the
fleet's own vehicles time; and making the signal commit to its timing costs the network more than the fleet recovers. A publication /
shadow comparison, an input intervention and an open-loop replay support a closed-loop account: once timing is published, the fleet's
slower approach changes the demand that the halting-based learned signal observes, the signal holds its greens longer, and
the forecast's own target shifts. We propose a two-sided interface (TSI) between the two frozen policies: a state-conditioned forecast for
the fleet, and a demand signal in front of the signal's inputs that also counts slow-moving vehicles. Neither policy is retrained; the
forecaster is fitted and the adapter tuned per deployed signal controller on training and tuning traffic (never on confirmation traffic). In the second registered confirmation
(the first, with one global adapter, failed on the fleet's time loss), on the same five signal checkpoints and four frozen fleets as the
first attempt with fresh traffic (family-wise bound 0.075),
TSI reduces the fleet's fuel by $3.85\%$ (95\% CI $[3.26, 4.43]$), city time loss by $0.55$\,s/veh and human time loss by $0.89$\,s/veh.
The fleet's time loss is non-inferior: $-0.22$\,s/veh, CI $[-0.44, +0.04]$, tested at one-sided $\alpha=0.025$ against a $+0.5$ margin.
Against a fleet-retraining recipe on the same panel, no fuel difference was detected (CI includes 0; this is not an equivalence test),
and all time losses were lower. In descriptive tests on grid4 with the design frozen, the limits were also met under a 25\% demand
increase, for rule and search-fitted fleets and at a 75\% fleet share; at a 25\% share the fleet's time-loss bound sits at the $+0.5$ margin
(non-inferiority not robustly shown). On a second network with multi-green learned signals and a rule fleet, the
adapter strongly reduces the harm of the forecast service, but the full configuration failed both registered attempts (time-loss limits
missed in both; no fuel reduction shown in the second). All results are simulated; no published method was reproduced under the same rights and budget.
\end{abstract}

\section{Introduction}
A connected-vehicle fleet trained to use signal phase and timing (SPaT) is useful only while the signals keep publishing timing it can act
on. Upgrading signals to learned adaptive controllers breaks this: the timing is no longer a static plan. Retraining the fleet each time a
city upgrades its signals, or retraining the signals each time a manufacturer updates its fleet, is often not possible; the two populations
are owned by different parties and released on different schedules. We ask whether a deployed learned fleet can keep its benefit when the
learned signals it relies on are upgraded, \emph{without retraining either population}.
\paragraph{Contributions.}
(i) A staged-upgrade benchmark with two learned populations (learned keep-or-switch signals; PPO and search-fitted
eco-driving fleets, with a fixed rule fleet as a non-learned receiver),
registered gates, and fresh-data confirmations (M1--M2).
(ii) A negative map of light-side and information-only repairs: commitments, gates, confidence and horizon filters, predictable splits and
per-decision commitment selection, each tested against the no-SPaT fallback; per-decision counterfactual credit was below the
simulator's chaotic noise floor in our measurement (signal-to-noise about 0.2).
(iii) A closed-loop diagnosis from three diagnostics (publication / shadow, an input intervention, an open-loop replay; no single mediator
isolated): once timing is published, the halting-based learned signal sees less of the fleet's demand and holds its greens longer, and the
forecast's own target shifts (a performative feedback); on a fixed schedule exact timing costs the fleet no time.
(iv) The two-sided interface TSI, confirmed in a second registered attempt on the same five light and four fleet checkpoints as the failed
first attempt, whose per-light results motivated per-light calibration (fresh traffic; family-wise bound 0.075), with component ablations, a same-panel comparison
with fleet retraining, and scope tests across receiver classes, penetration and demand.
\paragraph{Scope.} This is a methods paper with evidence on one network: the method failed on the second network, and no
published method was compared under the same rights and budget, so the results are neither general nor a state-of-the-art claim.

\section{Setting}
SUMO 1.27.1, default emission model; grid4 (16 signals, 1\,800 veh/h, 50\% fleet share, membership fixed per vehicle id); lights FT / ACT / learned (RL-H: a DQN that
keeps or switches between two greens from halting and vehicle counts, mean green about 15\,s against 42\,s under FT, C-100); fleets: SUMO Krauss at
imperfection $\sigma$ (non-learned family), ACC, and the learned fleets of M1 (R9, a four-parameter search policy; R10, a PPO speed
residual on the lane-aware GLOSA rule). Determinism and provenance: PYTHONHASHSEED = 0, bit-identical always-SUMO across jobs (C-27).
\paragraph{Seeds and panels.} Learned lights: development seeds k0 / k1, confirmation seeds t4--t8 (not used in TSI development; they trained the retrained
reference fleets F1-B (t4--t7) and served the first TSI attempt, whose per-light results motivated TSI-2; by design, TSI fits and tunes on
each of them, from TRAIN and TUNE traffic only).
Frozen fleets: development R10 s13 / s113, confirmation s213 / s313 / s413 / s513. M3 traffic: TRAIN 0--7 (fitting), DEV 600--611,
SELECT 650--655, TUNE 660--663 (per-light calibration), confirmation 700--719 (first attempt) and 720--739 (second attempt), demand shift
740--749 (period 1.6\,s), scope 750--759; every M3 traffic panel was generated after a non-existence check, and no traffic panel was used both
to choose and to confirm. The M1 / M2 panels (SELECT 100--104, fresh 450--509) are named where used.

\section{Milestone M1: a learned fleet worth deploying}
\paragraph{Criterion (registered).} Agent fuel superiority vs always-SUMO with agent and human time-loss margins (+0.5 s/veh),
teleports and unfinished trips; strong baseline $\sigma^*$ chosen on SELECT 100--104; confirmation on four new policy seeds and fresh traffic.
\paragraph{Fleet benefit (C-54, C-55; Table~\ref{tab:m1}).} Fresh traffic (20 seeds, generated after registration), FT, 50\% agents; learned
fleet = GLOSA-LD$\circ\sigma^*$, four policy parameters fitted by derivative-free search on training traffic (4 search seeds pooled).
Differences of the learned fleet are in Table~\ref{tab:m1}.
\begin{table}[t]\centering\small
\caption{M1: the learned fleet R9 (GLOSA-LD$\circ\sigma^*$) minus each comparator; fresh traffic (20 seeds), FT, 50\% agents; 95\%
intervals (C-54, C-55).}\label{tab:m1}
\begin{tabular}{lccc}\toprule
comparator & fuel (\%) & agent TL (s/veh) & human TL (s/veh)\\\midrule
always-SUMO & $-9.68$ [$-10.24$, $-9.00$] & $-3.48$ [$-3.77$, $-3.23$] & $-0.49$ [$-0.57$, $-0.41$]\\
$\sigma^*=0.3$ & $-7.58$ [$-8.14$, $-6.89$] & $-0.78$ [$-1.03$, $-0.58$] & $-0.36$ [$-0.45$, $-0.26$]\\
GLOSA (heuristic) & $-2.16$ [$-2.92$, $-1.35$] & $-3.03$ [$-3.88$, $-2.27$] & $-0.47$ [$-0.95$, $-0.07$]\\
GLOSA$\circ\sigma^*$ (no lane release) & $+0.16$ [$-0.57$, $+0.99$] & $-0.77$ [$-1.46$, $-0.16$] & $-0.30$ [$-0.65$, $+0.00$]\\
GLOSA-L$\circ\sigma^*$ (rule) & $-2.80$ [$-3.37$, $-2.06$] & $-0.18$ [$-0.41$, $+0.02$] & $-0.01$ [$-0.10$, $+0.08$]\\\bottomrule
\end{tabular}\end{table}
No teleports or unfinished trips for the learned fleet (GLOSA: 0.50, GLOSA$\circ\sigma^*$: 0.40 teleports per run). The fitted
parameters mainly learn where to release control (the last 12--22\,m before a needed lane change); the fuel gain over the rule is the
fuel the rule gives up by releasing early. Not RL; not compared with published eco-controllers; FT light only.
\paragraph{An RL fleet on the same panel (C-57).} A PPO speed residual on the rule (acting only where the rule commands; dense per-agent
cost; intervention-indexed discount, C-51), four new seeds pooled, against the rule: fuel $-1.70\%$ [$-2.33$, $-1.09$], agent time loss
$-0.04$ [$-0.29$, $+0.19$], human $+0.02$ [$-0.08$, $+0.12$]; against always-SUMO $-8.66\%$ / $-3.33$ / $-0.47$; against $\sigma^*$
$-6.54\%$ / $-0.64$ / $-0.33$; no teleports. Against GLOSA$\circ\sigma^*$ without lane release it is $+1.29\%$ fuel: it does not recover the fuel the
rule gives up by releasing early, which R9's fitted release distance does.
\paragraph{SPaT ablation.} For the fitted GLOSA class, removing SPaT means the policy never commands and reduces to the
$\sigma^*$ fleet; the comparison against $\sigma^*$ above is therefore the class-level SPaT ablation ($-7.6\%$ fuel). It does not separate
information from representation. The PPO residual (R10, confirmed in M1, C-57) reduces to the $\sigma^*$ fleet exactly when no SPaT is
published (bit-identical rows, M2 check KM2-2), so its benefit against $\sigma^*$ is likewise the class-level SPaT ablation.
\paragraph{An actuation pitfall shared by every commanded fleet (C-45).} With SUMO's \texttt{setSpeed}, every commanded fleet we
trained or wrote (GLOSA, an in-class controller, PPO, behaviour clones, earlier learned fleets) created teleports (0--17 per run on TRAIN 0--7;
GLOSA 0--2, PPO snapshots 4--14), while always-SUMO and $\sigma=0.3$ created none. Each inspected event is a lane-change deadlock: an agent car stopped at the end of a lane that does
not continue on its route. Releasing control whenever the car's current lane does not continue (a code predicate on SUMO's best-lanes
offset) removed these deadlocks on every panel tested (C-46). This is trace and intervention evidence in this SUMO version; without a
same-checkpoint factorial over the two interfaces we do not attribute the earlier learned fleets' time-loss failures to it in full (C-50).
\paragraph{A strong same-information rule (C-46).} GLOSA on the $\sigma^*$ base driver with lane-aware release (command only on a lane
that continues) is beneficial on fresh traffic (20 seeds): against $\sigma^*$, agent fuel $-5.0\%$ and agent / human time loss $-0.7$ /
$-0.4$\,s/veh; against always-SUMO, $-7.1\%$ and $-3.3$ / $-0.5$; no teleports. It is a fixed rule, not learned, not compared with
published eco-controllers, and without a SPaT ablation of its own it does not show that SPaT is necessary; it is the baseline a learned
fleet must beat on the same base and interface.
The learned four-parameter policy (R9) against this rule on fresh traffic is in Table~\ref{tab:m1} (C-54).

\section{Milestone M2: is there a compatibility problem?}
\paragraph{Design.} The registered target was a $2\times2$ matrix (old / upgraded light, frozen / adapted fleet) with both lights learned and
beneficial in their reference condition, absolute effects and difference-in-differences. The full matrix was not run: M2a and M2b below are
its parts that were run, the adapted-fleet column is the retrained reference F1-B in M3, and a learned old light $L_0$ was not run (C-60).
\paragraph{M2a: SPaT service withdrawn or misused (C-59, C-60; fresh traffic, 20 seeds; four light seeds paired with four fleet
seeds, not a factorial).} Scope of this panel: the old light is fixed-time (not a learned L0); it has no forecast arm (run later as M2b)
and no adapted-fleet column (run later as the retrained reference in M3). The learned light is a large
upgrade over the fixed-time program with SUMO cars (city time loss $-33.7$\,s/veh, city fuel $-19.5\%$). Both confirmed fleets lose their
own benefit. If the learned light correctly publishes no plan, each fleet falls back to its base driver and becomes identical to
$\sigma^*$: the benefit drops from $-7.5\%$ (fitted fleet) and $-6.5\%$ (PPO fleet) agent fuel to exactly 0, with no harm. If the fleet
keeps deriving SPaT from the static program, which the learned light no longer follows (30--43\% of its SPaT calls are marked valid
outside the function's guaranteed scope; these are calls, not measured forecast errors), it is harmed against $\sigma^*$ under the same light: $+7.4\%$ / $+6.5\%$ fuel and $+1.7$ / $+1.6$\,s/veh agent time loss. In
absolute terms the agents still gain from the upgrade; what is lost is the fleet's benefit against its counterfactual drivers.
\paragraph{M2b: an empirical forecast from the learned light (C-62, C-71, C-72).} A conditional-quantile forecast of the learned light's
green durations does not return the benefit to the frozen fleet; at the median it harms both fleets against $\sigma^*$ (R9 $+7.12\%$ fuel and
$+3.64$\,s/veh agent time; R10 $+4.32\%$ and $+3.37$; fresh 490--509, C-62). Two facts are consistent with this: the original forecast was 1\,s early in
the green branch on 100\% of points (a timing-convention error, fixed in a corrected provider that still left one yellow branch 1\,s early, C-84; the provider c2 used from
the retrained reference F1-B on, and by TSI, corrects both),
and the frozen fleet is sensitive to timing errors, with asymmetric fuel and time effects in an open-loop oracle-shift diagnostic (exact
timing $-6.6\%$ fuel; 2\,s early $-2.4\%$; 2\,s late $-4.2\%$ fuel but agent time loss $+1.5$\,s/veh; 8\,s late $+9.3\%$ and $+11.6$\,s/veh;
C-72). The diagnostic uses another fleet and schedule than the retrained fleet, so it warns of the trade-off but does not predict it.
\paragraph{Second network (C-81).} On an independent own-built network (SUMO \texttt{netgenerate} spider, 18 signals, six with three greens),
upgrading the static programs to a learned all-green controller cuts city time loss by 58.9\,s/veh. The frozen rule fleet (GLOSA-L$\circ\sigma^*$,
$-9.2\%$ agent fuel under the static programs) again loses its benefit: labelled $\to$ exactly $\sigma^*$ (no harm); naive $\to$ harmed
($+7.2\%$ fuel, agent TL $+2.7$, human TL $+0.6$, upper $+0.8$).
\begin{figure}[t]\centering\includegraphics[width=\linewidth]{figs_two_pop/fig1_m2_crossplay.pdf}
\caption{The frozen fleet's benefit against $\sigma^*$ under the same light, before (FT) and after the learned-light upgrade. Labelled bars are
exactly 0 (rows identical to $\sigma^*$). 95\% intervals; grid4 fresh 470--489: pair $\times$ traffic bootstrap over the four paired light / fleet
seeds (not crossed); spider fresh 1000--1019: crossed over two lights $\times$ traffic. Script \texttt{scripts/CITY\_fig\_m2\_crossplay.py}.}\end{figure}
\section{Milestone M3: recovering part of the frozen fleet's benefit without retraining either population}
\paragraph{Endpoint (registered).} Agent fuel against R, the upgraded learned light publishing no SPaT (the frozen fleet then drives as
$\sigma^*$), with city, agent and human time-loss non-inferiority at $+0.5$\,s/veh (one-sided upper bounds), teleports and unfinished
trips; crossed bootstrap over light seed, fleet seed and traffic seed.

\subsection{What does not work (negative map)}
\paragraph{Light-side commitments.} Reward-shaped lights, fixed-window and ETA-triggered commitments, a predictable cycle and
receiver-triggered precision commitment (RPC) failed their DEV gates; a registered negative confirmation on fresh traffic found WC5 and RPC worse than the fallback
($+2.98\%$ and $+9.26\%$ agent fuel, C-86). A $2\times2$ diagnostic of RPC (C-95) shows both suspected costs are real (the forced hold:
$-4.54\%$ fuel when removed; withdrawing the message at yellow onset: $-2.80\%$ when kept, $-2.66\%$ with the light schedule fixed), yet the best
cell is still $+3.06\%$: the commitment's own cost exceeds what the fleet recovers.
\paragraph{Choosing commitments decision by decision.} Re-simulating full episodes from identical prefixes with one intervention
(96 situations, five options, C-96) shows that one 5--15\,s commitment changes 640--900 of 1\,800 vehicles. Out of sample, no selection
rule (sensitivity, cost, both, or a learned regressor) was better than a random option (every interval includes 0), and the best fixed option does not beat a permutation null
($p=0.29$; within-situation spread $0.235\%$ vs $0.040\%$ between options, C-98). Single-decision gains also do not add up: the
additive prediction from single edits ($-32$ to $-176\%$, intervals of about $\pm200\%$) carries no information about RPC's realized
$+9.4\%$ on the development panel (post hoc, C-101). Candidates must therefore be evaluated as
whole policies over full episodes.
\paragraph{Information-only repairs.} Abstaining where the forecast leaves its support (0.11\% of the fleet's queries, C-97), confidence and
horizon gates (every gate moves along one fuel-for-agent-time frontier, C-104) and a predictable split controller tuned by CEM (it gives the
fleet $-7.9\%$ fuel from exact timing, but the light alone costs $+2.1$\,s/veh city time loss against the learned light, C-116); none
reached a feasible point in the families and budgets tried.

\subsection{Diagnosis: a closed-loop account}
A state-conditioned direct time-to-green forecaster (an MLP quantile model on phase, age, queue and recent green history, fitted on training
traffic, C-99) lowers the median-forecast error by a third (the 0.1-quantile error is not lower, so the
registered accuracy criterion failed, C-99) and gives the frozen fleet $-3.4\%$ fuel against R, but costs its vehicles
$+2.3$\,s/veh while human time falls ($-1.3$). Three diagnostics support a closed-loop account; they do not isolate a single mediator,
since the adapter used in (a) and (c) changes the light's behaviour more broadly than demand perception alone.
(a) With SPaT published, mean halting on the served lanes falls from 2.9 to 1.6 and the learned light's greens lengthen (14.8 to 16.9\,s; C-99),
consistent with eco-approaching vehicles moving above the $0.1$\,m/s halting threshold that the light observes and was trained on. Counting slow vehicles
(0.1--2\,m/s) as waiting in front of the frozen light cuts the agents' time cost by $0.8$\,s/veh (C-100).
(b) On the light's recorded schedule replayed open loop, exact timing gives the fleet $-5.6\%$ fuel \emph{and} $-0.37$\,s/veh, while the forecaster
costs $+2.7$\,s/veh on the same schedule: on that schedule the time cost comes from the forecast's errors, not from a time-for-fuel
trade-off of the fleet (C-100, C-107).
(c) Publishing the forecast shifts its own target: at the fleet's queries the error is $3.02$\,s when published vs $1.70$\,s when computed
but not used, and scored waits longer than 30\,s rise from $0.09\%$ to $3.34\%$ of scored queries (C-109, C-113); the demand adapter (DEV, global $q{:}2{:}0$) roughly halves this
sensitivity ($+0.73$ vs $+1.32$\,s, difference $-0.59$\,s $[-0.88,-0.35]$, C-122; intervention evidence, not a measured loop gain or a stability proof). Across the five confirmation lights of the first attempt, the residual time cost grows with the residual green
lengthening (five points, descriptive; Figure~\ref{fig:mech}a, C-117).

\subsection{Method: a two-sided interface (TSI)}
Both learned populations keep their weights. (1) \emph{Fleet side}: a hybrid forecast, the program's exact answer where the schedule is
determined (the green after a yellow) and the state-conditioned forecaster elsewhere, served at the median. (2) \emph{Light side}: a demand
adapter in front of the frozen light's halting inputs, $q = h_{<0.1} + A\,n_{[0.1,2)} + B\,n_{[2,5)}$ on the served lanes. (3)
\emph{Calibration per deployed light from training and tuning traffic only}: the forecaster is fitted on eight training episodes of that light; $(A,B)$
is chosen on four tuning episodes by a pre-registered rule: among the grid points whose agent time-loss point is $\le -0.25$\,s/veh and
whose city and human points are $\le 0$ against R, the most fuel saving; otherwise the smallest agent time loss. Calibration uses the
deployed light itself (that is the method), a development fleet (R10 s13) and TRAIN / TUNE traffic; no confirmation traffic and no
confirmation fleet enters it.

\paragraph{Calibration cost (per deployed light controller).} Forecaster: 8 training episodes with the fleet using the deployed age-only
forecast (feature log), one fit of the 64--64 MLP (CPU, about 30\,s; 40 epochs). Adapter: a grid of 10 $(A,B)$
pairs on 4 tuning episodes plus R (44 episodes). No policy gradient step for either population. Sensors: halting and vehicle counts already
used by the light, plus vehicle speeds on the served lanes (three speed bins). Inference: one MLP pass per light and second when the fleet
queries; the exact program answer in the determined case. Implementation note: the grid4 provider assumes two greens per light; the spider
version generalizes the features to two or three greens (Section~M4).

\subsection{Confirmation}
A first registered attempt with one global adapter chosen on two development lights (C-110) failed on agent time only: on five lights and
four fleets not used in TSI development and fresh traffic (700--719), fuel $-4.56\%$, city $-0.20$, human $-1.23$, agent $+0.85$\,s/veh (upper $+1.27$, C-114). The
second and last registered attempt (TSI as above, per-light calibration; traffic 720--739; $\alpha=0.025$ one-sided because it is a second
attempt, Bonferroni bound $0.075$ over both) passed every limit (Table~\ref{tab:tsi}, Figure~\ref{fig:tsi}, C-118).
\begin{table}[t]\centering\small
\caption{TSI confirmation (C-118), the second registered attempt (the first, with one global adapter, failed on agent time, C-114). Both
attempts use the same five learned-light checkpoints and four frozen fleets (fresh traffic each time); $\alpha=0.025$ one-sided here,
Bonferroni bound $0.075$ over both. Five learned-light seeds $\times$ the four frozen PPO fleets s213--s513 (R10's M1 confirmation seeds)
$\times$ 20 fresh traffic seeds, crossed
95\% intervals; the primary test uses the upper end (one-sided 97.5\%). F1-B: the fleet retrained under the new light (more update rights),
same panel; it is the repository's retraining recipe (provider and pipeline differ), not every retraining method. TSI $-$ F1-B on fuel:
no difference detected, not an equivalence test.}\label{tab:tsi}
\resizebox{\linewidth}{!}{\begin{tabular}{lcccc}\toprule
vs R & agent fuel (\%) & city TL (s/veh) & agent TL (s/veh) & human TL (s/veh)\\\midrule
TSI & $-3.85$ [$-4.43$, $-3.26$] & $-0.55$ [$-0.67$, $-0.44$] & $-0.22$ [$-0.44$, $+0.04$] & $-0.89$ [$-1.10$, $-0.67$]\\
service only (HYB) & $-4.93$ & $+0.26$ & $+1.97$ & $-1.43$\\
adapter only (ADP) & $+0.29$ & $-0.17$ & $-0.09$ & $-0.25$\\
fleet retrained (F1-B) & $-3.18$ & $-0.17$ & $+0.28$ & $-0.62$\\\midrule
TSI $-$ F1-B & $-0.68$ [$-1.70$, $+0.11$] & $-0.38$ [$-0.59$, $-0.18$] & $-0.50$ [$-0.88$, $-0.09$] & $-0.27$ [$-0.44$, $-0.06$]\\\bottomrule
\end{tabular}}\end{table}
\begin{figure}[t]\centering\includegraphics[width=\linewidth]{figs_two_pop/fig3_tsi_components.pdf}
\caption{TSI against the fallback R on the confirmation panel, with its two halves and the retrained fleet (95\% intervals; dashed: the
$+0.5$\,s/veh margin). Script \texttt{scripts/CITY\_fig\_tsi.py}.}\label{fig:tsi}\end{figure}
In the ablations tried the two sides are complementary, and neither alone meets the joint objective: the fleet-side service carries the fuel ($-4.12\%$ against the adapter alone) and the light-side adapter carries
the agents' time ($-2.19$\,s/veh against the service alone). The per-light chosen adapters are $(1.5,0.25)$, $(1,0.5)$, $(1,0.5)$,
$(1,0.25)$, $(1.5,0.5)$; agent time against R is within $[-0.47,+0.18]$\,s/veh on every light.
\begin{figure}[t]\centering\includegraphics[width=\linewidth]{figs_two_pop/fig4_tsi_mechanism.pdf}
\caption{Mechanism. (a) First TSI attempt (global adapter) and the service alone: per light, the agents' time cost grows with how much the
learned light lengthens its greens (C-117). (b) On a fixed schedule exact timing saves fuel and time; the forecast's errors cost time, and in
closed and open loop differ by $-3.35\%$ agent fuel (C-100, C-107).}\label{fig:mech}\end{figure}

\subsection{Scope}
With the TSI design frozen (Table~\ref{tab:allstats}, Figure~\ref{fig:scope}): a 25\% demand increase on fresh traffic, $-3.79\%$ fuel and $-0.27$\,s/veh agent
time (C-120); a rule fleet, $-1.84\%$ and $+0.03$; a search-fitted fleet, $-1.38\%$ and $+0.13$ (3 teleports in 100 runs against 0 for R, within the $+0.1$ per-run limit; the GEN report's ``all limits'' line did not check teleports, which were counted afterwards from the saved rows, C-126); fleet shares of 25 and 75\%, $-4.13\%$
/ $+0.07$ (one-sided upper $+0.4986$ with the registered bootstrap draws; re-drawn with 40 bootstrap seeds the upper is
$+0.51\pm0.01$ and is $\ge+0.5$ for 35 of 40, so non-inferiority at the 25\% share is at the margin and not robustly shown; two-sided
interval $[-0.36, +0.59]$; C-128) and $-4.75\%$ / $-0.58$ (C-121). Replacing the state forecaster by the age-only forecast with the same
adapter harms ($+6.72\%$; TSI is $9.87$ points better): the age-only forecaster tried does not replace the state forecaster in this
design (not a claim about every forecaster without state). The reverse upgrade order met its registered no-harm criterion
in this setting: a fleet retrained for the new light, under the old light and against the original R10
fleets, $-2.23\%$ fuel (upper $-1.62$), agent time $+0.15$ (upper $+0.40$; C-121).
\begin{table}[t]\centering\small
\caption{Every TSI test against R, with one convention. Fuel: agent fuel change (\%), point (upper bound); time losses: upper bound
(s/veh); the limits are fuel upper $<0$, every time-loss upper $<+0.5$, teleports $\le$ R $+0.1$ per run. Upper bounds are at the test's
level: one-sided $\alpha=0.05$ (95\%) or $0.025$ (97.5\%) for the registered confirmations, one-sided 95\% for the descriptive scope
tests. Crossed bootstrap over light seed, fleet seed and traffic seed; panel = lights $\times$ fleets $\times$ traffic seeds. Every
grid4 panel reuses the light checkpoints t4--t8 and (except R8, R9) the PPO fleets s213--s513, so the rows are not independent
replications. The spider attempt-2 report header prints the traffic range of attempt 1; the data are 1130--1149 (C-124).}\label{tab:allstats}
\resizebox{\linewidth}{!}{\begin{tabular}{llllrrrrl}\toprule
test & panel & level & agent fuel & city & agent & human & teleports/run (R) & verdict\\\midrule
TSI, attempt 1 (C-114) & grid4, $5\times4\times20$ (700--719) & reg., 0.05 & $-4.56$ ($-4.30$) & $-0.05$ & $+1.27$ & $-1.06$ & 0 (0) & FAIL (agent)\\
TSI-2 (C-118) & grid4, $5\times4\times20$ (720--739) & reg., 0.025 & $-3.85$ ($-3.26$) & $-0.44$ & $+0.04$ & $-0.67$ & 0 (0) & PASS\\\midrule
demand $+25\%$ (C-120) & grid4, $5\times4\times10$ (740--749) & desc., 0.05 & $-3.79$ ($-3.26$) & $-0.52$ & $-0.07$ & $-0.86$ & 0 (0) & within limits\\
rule fleet R8 (C-121) & grid4, $5\times1\times10$ (750--759) & desc., 0.05 & $-1.84$ ($-1.50$) & $-0.28$ & $+0.24$ & $-0.64$ & 0 (0) & within limits\\
search fleet R9 (C-121) & grid4, $5\times2\times10$ (750--759) & desc., 0.05 & $-1.38$ ($-0.58$) & $-0.28$ & $+0.36$ & $-0.73$ & 0.03 (0) & within limits\\
fleet share 25\% (C-121) & grid4, $5\times2\times10$ (750--759) & desc., 0.05 & $-4.13$ ($-3.78$) & $-0.20$ & $+0.4986$ & $-0.30$ & 0 (0) & at the margin, not robust (C-128)\\
fleet share 75\% (C-121) & grid4, $5\times2\times10$ (750--759) & desc., 0.05 & $-4.75$ ($-4.43$) & $-0.47$ & $-0.33$ & $-0.67$ & 0 (0) & within limits\\\midrule
spider, attempt 1 (C-123) & spider, $2\times1\times20$ (1110--1129) & reg., 0.05 & $-1.66$ ($-1.43$) & $+0.91$ & $+1.94$ & $+0.01$ & 0.025 (0) & FAIL (city, agent)\\
spider, attempt 2 (C-124) & spider, $2\times1\times20$ (1130--1149) & reg., 0.025 & $-0.31$ ($+0.16$) & $+1.77$ & $+1.94$ & $+1.64$ & 0 (0) & FAIL (all four)\\\bottomrule
\end{tabular}}\end{table}
\begin{figure}[t]\centering\includegraphics[width=\linewidth]{figs_two_pop/fig5_tsi_scope.pdf}
\caption{TSI with its design frozen across receiver classes, fleet shares and a demand shift (fresh traffic; 95\% intervals).}\label{fig:scope}\end{figure}

\section{Milestone M4: a second setting}
Spider network (own-built, \texttt{netgenerate}): a learned all-green light beats the static programs by 57\,s/veh city time loss (two light seeds, development
traffic, C-76); the rule fleet saves $9.2\%$ fuel against $\sigma^*$ (development traffic, C-77), with agent / human time-loss non-inferiority not established on fresh traffic
(upper $+0.55$ / $+0.52$ against the $+0.5$ margin, C-81);
the four-parameter search fleet (R9) failed its pilot gate against the rule there (fuel $+0.40\%$ $[+0.25,+0.53]$ on development
traffic, C-80); the M2 loss and harm replicate with the rule fleet (C-81); fixed-window commitment is again worse than the
no-SPaT fallback, with a different decomposition (no fleet contribution detected, C-82). The retrained-fleet check was not run there, because its registered condition (F1-A1 confirmed; M4-A7) failed (C-85). The
earlier RESCO cologne8 attempt failed at the light baseline and is kept as a negative (C-35, C-36).
\paragraph{TSI on the second network (C-123, C-124).} The same procedure (the forecaster generalized to two or three greens, the adapter
in front of the multi-green learned light's current / next-green inputs, per-light calibration on new tuning traffic) was tested in two registered
confirmation attempts on fresh spider traffic with the rule fleet; with two light seeds, both are conditional on these two lights. The service alone is strongly harmful there (attempt 1: $+20.7\%$ fuel, $+59.5$\,s/veh agent time, with about one
teleport per run; attempt 2: $+19.8\%$, $+57.8$); adding the adapter reduces that harm (TSI $-$ service: $-58.1$ / $-56.4$\,s/veh agent time), but the full configuration failed both attempts. Attempt 1: fuel $-1.66\%$
$[-1.99,-1.38]$, and both city time loss ($+0.55$, upper $+0.91$) and agent time loss ($+1.33$, upper $+1.94$) fail; $0.025$ teleports per
run (R: 0). Attempt 2, with a wider adapter grid at $\alpha=0.025$: fuel $-0.31\%$ $[-0.81,+0.16]$ (no fuel reduction shown); city $+1.32$,
agent $+1.43$ and human $+1.20$\,s/veh all fail. Spider changes the network, the phase structure, the light learner and the receiver at
once, so the failure is not attributed to any one of them. The reduction of harm is relative to a very harmful service; it is not a
deployable result.

\section{Related work}
Joint or cooperative signal--vehicle learning trains both populations together: CoTV [7] controls signals and the nearest connected
vehicle on each approach with multi-agent deep RL, and HONEST-CAV [11] optimizes network signals and vehicle trajectories hierarchically
with value-decomposition MARL. Our setting differs in its rights, not its ambition: both populations are frozen and upgraded on
separate schedules. SPaT prediction for actuated signals is established, with supervised time-to-green models [6] and deep ensembles that
attach a confidence level [5]; our forecaster is not new and meets the joint objective only in combination. Observation adapters in front of frozen
policies also exist: DC-Ada adapts per-robot observation transforms for a frozen shared policy by reward-only random search [1], and
adaptive modular world models separate a site-specific observation adapter from a shared dynamics model for traffic-signal transfer [8].
Ours is a two-parameter physical count transform calibrated against a partner-induced shift rather than a sensing-pipeline change. That
published forecasts change their own target is the consistency problem of anticipatory route guidance [2] and of performative prediction
[10, 4]; we identify a traffic-signal instance and show, on development traffic, that a light-side input change reduces the measured
sensitivity. The complementarity of
the two interfaces is consistent with a small-gain reading of the loop forecast $\to$ fleet $\to$ light $\to$ target [9, 3]; this is
motivation, not a theorem about the learned light. No published method was reproduced under the same rights and budget, so we make no
state-of-the-art or ``first'' claim; forecaster + adapter as a combination is not by itself novel, and the contribution is the
closed-loop diagnosis and the confirmed two-sided calibration procedure in this setting.

\section{Limitations}
Simulation only (SUMO 1.27.1, default emission model). The method is confirmed on one network (grid4); on the second network (spider, multi-green
lights, rule fleet) the full configuration failed both registered attempts (time-loss limits in both; no fuel reduction shown in the second;
C-123, C-124). The confirmation that
passed is a second attempt; its light and fleet checkpoints are those of the failed first attempt (the traffic is fresh), and the
family-wise bound over the two attempts is $0.075$. The crossed bootstrap rests on five light and four fleet checkpoints. No published
method was run under the same rights and budget. Against the repository's retraining recipe, fuel equivalence is not shown. Calibration needs 52 simulated episodes of each deployed light with a development fleet, so a deployment needs a site simulator; the
transfer of this calibration from a simulator to a real site is untested. The demand adapter needs per-lane speed sensing at the light. One learned-light family per network (keep/switch DQN on halting counts on grid4); a
light trained on other inputs may not shift the same way. The provider of M2b and of the first retrained-fleet attempt leaves a 1\,s early yellow branch (C-84;
corrected in the provider c2 used by F1-B and TSI); the hybrid is exact only for the green after a yellow (C-113). Per-decision analyses are limited by the simulator's chaotic divergence
(C-98).

\section{Reproducibility}
Registered endpoints come from registered report scripts over saved per-episode rows; every other number comes from the same saved rows
through the analysis named in its ledger row, and the post-hoc ones are labelled (C-55, C-101, C-117, C-125, C-126, C-128).
Provenance gap: the TSI, TSI-2, SHIFT, GEN and spider protocols have no launch-time hash record, and the TSI / HYB rows do not carry the
forecaster, fleet or table hashes (the evaluator's metadata is not written when it is imported; C-128). Every recorded component hash equals
the current source, and every module and input on these paths predates its first use (file times); the freeze record of C-128 lists them. Main files (repository paths):
\begin{itemize}\setlength\itemsep{0pt}
\item protocols: \path{research/city_sim/PROTOCOL_CITY_TSI2_20261005.md} (method and confirmation) and the TSI, SHIFT, GEN, PERF and
SPIDER\_TSI protocols of the same date in the same directory;
\item evaluators: \path{scripts/CITY_eval_n4.py} (grid4 service and \path{--light-obs q:A:B}), run through
\path{scripts/CITY_f1a1_run.py} (loads the light checkpoints t4--t8) and, for fleet shares, \path{scripts/CITY_rho_run.py}; \path{scripts/CITY_m4s_tsi_eval.py} (spider);
\item forecasters: \path{scripts/CITY_n2_fit.py}, \path{scripts/CITY_n2k.py}; calibration: \path{scripts/CITY_tsi2_tune.py};
\item reports: \path{scripts/CITY_tsi2_report.py}, \path{CITY_shift_report.py}, \path{CITY_gen_report.py}, \path{CITY_perf_report.py},
\path{CITY_sptsi_report.py}; figures: \path{scripts/CITY_fig_tsi.py}.
\end{itemize}
Determinism: \path{PYTHONHASHSEED=0}, \path{OMP_NUM_THREADS=1}. Parity checks KN2 / KT-1 / KS-1 compare offline and online features and
predictions (0 mismatches).

\section*{References}
\small
\begin{enumerate}\setlength\itemsep{1pt}
\renewcommand{\labelenumi}{[\arabic{enumi}]}
\item S.~Alqithami. DC-Ada: Reward-only decentralized sensor adaptation for heterogeneous multi-robot teams. \emph{IEEE Access}, 2026.
doi:10.1109/ACCESS.2026.3687163; arXiv:2604.03905.
\item J.~A.~Bottom. \emph{Consistent Anticipatory Route Guidance}. PhD thesis, Massachusetts Institute of Technology, 2000.
\item S.~Dashkovskiy, B.~S.~R\"uffer, F.~R.~Wirth. An ISS small gain theorem for general networks. \emph{Math. Control Signals Systems}
19(2):93--122, 2007.
\item D.~Drusvyatskiy, L.~Xiao. Stochastic optimization with decision-dependent distributions. \emph{Math. Oper. Res.} 48(2):954--998, 2023.
\item S.~Eteifa, A.~Shafik, H.~Eldardiry, H.~A.~Rakha. Deep learning ensemble approach for predicting expected and confidence levels of
signal phase and timing information at actuated traffic signals. \emph{Sensors} 25:1664, 2025. doi:10.3390/s25061664.
\item A.~Genser, M.~A.~Makridis, K.~Yang, L.~Amb\"uhl, M.~Menendez, A.~Kouvelas. Time-to-green predictions for fully-actuated signal
control systems with supervised learning. \emph{IEEE Trans. Intell. Transp. Syst.}, doi:10.1109/TITS.2023.3348634; arXiv:2208.11344.
\item J.~Guo, L.~Cheng, S.~Wang. CoTV: Cooperative control for traffic light signals and connected autonomous vehicles using deep
reinforcement learning. \emph{IEEE Trans. Intell. Transp. Syst.} 24(10), 2023. arXiv:2201.13143.
\item Z.~Huang, Y.~Liu, C.~Liang, G.~Zheng. Planning under observation mismatch for traffic signal control via adaptive modular world
models. \emph{ICAPS}, 2026. arXiv:2501.02548.
\item Z.-P.~Jiang, A.~R.~Teel, L.~Praly. Small-gain theorem for ISS systems and applications. \emph{Math. Control Signals Systems}
7(2):95--120, 1994.
\item J.~Perdomo, T.~Zrnic, C.~Mendler-D\"unner, M.~Hardt. Performative prediction. \emph{Proc. ICML}, PMLR 119:7599--7609, 2020.
\item Z.~Zhang, C.~Wan, P.~Hao, K.~Boriboonsomsin, M.~J.~Barth, Y.~Liu, S.~Ucar, G.~Wu. HONEST-CAV: Hierarchical optimization of network
signals and trajectories for connected and automated vehicles with multi-agent reinforcement learning. arXiv:2602.18740, 2026.
\end{enumerate}
\normalsize

\appendix
\section{History of the M1 and M2 recipes}
\paragraph{Earlier fleet recipes (C-21 \dots C-37).} No earlier recipe was confirmed (C-21 \dots C-31, corrected C-37); one development checkpoint (R5h\_t0)
passed on fresh traffic; the non-learned $\sigma$-family is a strong baseline ($\sigma=0.3$: $-2.5$ s/veh, $-2.2\%$ fuel).
\paragraph{Fleet diagnostics (C-39 \dots C-43).} D1: the ported phase reward (B) prefers physically better behaviour for large
differences but does not encode the time-loss constraints (it ranks an infeasible fuel-for-time policy far above always-SUMO); PPO from
scratch did not reach always-SUMO even on its training traffic. D2: a hand-written in-class controller (GLOSA target, no release) gives
$-6.7\%$ fuel at equal agent time but harms human cars; a behaviour clone of the heuristic GLOSA trades fuel for time in closed loop.
An evaluation bug (exit release missing in deterministic evaluation) was found by an identity check and fixed (C-41).
A residual on the in-class controller with a crawl-proof group cost (v9) failed its pilot gate; its confirmation was not evaluated (C-47).
Derivative-free policy search over a three-parameter GLOSA family on the deployment objective (fuel with time-loss penalties) returned
the untuned heuristic itself for the pilot search seed; the heuristic fails the $\sigma^*$ comparison on agent time loss (+1.8\,s/veh,
SELECT) and the human-time-loss margin on the width of its 5-seed bound (C-43).
\paragraph{Earlier compatibility probes (C-14 \dots C-34).} Weak learned fleets regressed under light upgrades and a car-side reaction repair
helped (C-14, C-16); non-learned fleet upgrades ($\sigma$, ACC) were not harmed by three learned-light families (C-30, C-32 \dots C-34).

\section{Candidates that failed, and errors found and fixed}
\paragraph{Failed candidates (all reported in the ledger).} M3-E reward-shaped lights (DEV); fixed-window, ETA and receiver-triggered
commitments (WC5, RPC confirmed worse than the fallback, C-86; RPC $2\times2$ diagnostic C-95); per-decision commitment selection (no
selection rule better than a random option out of sample, every interval includes 0; C-96, C-98); tail-support abstention (C-97);
confidence / horizon gates and gate + adapter (C-104, C-105); a CEM-tuned
committed-duration controller (C-116); TSI variants (direct-only hybrid, no SPaT on permissive lefts; C-119); TSI first confirmation
(C-114); spider attempts 1 and 2 (C-123, C-124).
\paragraph{Errors found and fixed.} Found after results that used them, and disclosed with those results: the SPaT timing convention
($+1$\,s, C-71; M2b used the earlier provider) and a yellow branch left 1\,s early by the first correction (C-84, found by an audit;
corrected in the provider c2 for all later runs). Found by self-checks or job failures before the verdict they affected: an xargs job-merging bug (C-92); an exactness error in the committed-duration controller's SPaT (two fixes, KN8-1, C-112);
missing output directories (C-112); a tuple error in the spider provider, caught by KS-2 (C-123 Amendment A1). One provider file was edited
while a job imported it, with no change to the registered code paths (C-103).
Reporting errors corrected in the text: ``exact in yellow'' had been claimed from the 0.1-quantile table (C-113); the spider attempt-2
report header prints attempt 1's traffic range (C-124); the scope report's limits line skipped teleports (C-126); the demand-shift report tested neither teleports nor unfinished trips (both 0
in its 600 runs, C-128); the spider attempt-2 amendment states it was registered before attempt 1's result was read, which the file
times cannot verify (attempt 1's per-seed log lines were printed before the attempt-2 job list was written; C-128).
\section{Validation and pitfalls (QA)}
Reward traces, sentinel readings, residual ratchet, horizon semantics, dominance labels (C-17, C-22, C-25, C-27, C-37, C-38).

\end{document}
